% --- Archon physics formalization source begin ---
% archon:physics
% archon:covers PhyXMiniProblems/problem_phyx_mini_0291.lean
% archon:source-report reports/phyx_mini/problem_phyx_mini_0291.source.json

\chapter{Physics problem phyx\_mini\_0291}
\label{ch:PhyXMiniProblems_problem_phyx_mini_0291}

\paragraph{Problem source.}
\#\# Physical scenario

A sinusoidal transverse wave is traveling along a string in the negative direction of an \$x\$ axis. Figure shows a plot of the displacement as a function of position at time \$t = 0\$; the scale of the \$y\$ axis is set by \$y\_s = 4.0\$ cm. The string tension is \$3.6\$ N, and its linear density is \$25\$ g/m.

\#\# Question

Find the period of the wave.

\#\# Answer choices

- A: A: 0.019 s

- B: B: 0.028 s

- C: C: 0.045 s

- D: D: 0.033 s

\#\# Figure caption (auxiliary; use the image as primary evidence)

The image is a graph depicting a sinusoidal wave. Here are the details:

- **Axes:** 
  - The horizontal axis represents \textbackslash{}(x\textbackslash{}) in centimeters (cm), labeled as \textbackslash{}(x \textbackslash{}, (\textbackslash{}text\{cm\})\textbackslash{}).
  - The vertical axis represents \textbackslash{}(y\textbackslash{}) in centimeters (cm), labeled as \textbackslash{}(y \textbackslash{}, (\textbackslash{}text\{cm\})\textbackslash{}).

- **Grid:** The background features a grid to help identify the coordinates.

- **Waveform:** 
  - The wave is sinusoidal, smoothly oscillating above and below the horizontal axis.
  - The wave appears to have an amplitude labeled as \textbackslash{}(y\_s\textbackslash{}) and \textbackslash{}(-y\_s\textbackslash{}), indicating the maximum and minimum heights of the wave concerning the horizontal axis.

- **Plot Range:**
  - The wave completes roughly one full cycle between \textbackslash{}(x = 0\textbackslash{}) and \textbackslash{}(x = 40\textbackslash{}) cm on the horizontal axis.

- **Labels:**
  - Specific points along the horizontal axis are marked, such as 0, 20, and 40, likely indicating significant points for the wave's cycle.

The combination of these elements suggests a typical sinusoidal wave, as might be found in physics or engineering contexts, representing phenomena such as waves, oscillations, or alternating current signals.

\#\# Dataset metadata

Wave Properties; Physical Model Grounding Reasoning, Multi-Formula Reasoning

\paragraph{Recorded answer/context.}
D: D: 0.033 s

\paragraph{Figure/image path.}
/root/proposal\_for\_physic/science-mango/phyx\_mini\_run/phyx\_data/test\_image/291.png

\paragraph{Formalization target.}
create a compiling Lean file with sorry bodies at `PhyXMiniProblems/problem\_phyx\_mini\_0291.lean`. The Lean declarations must preserve the physical quantities, dimensions or dimensional roles, figure labels, governing-law hypotheses, and final relation expressed by this problem.
Use Mathlib/Physlib names found through LeanExplore where available. If a physics API is missing, introduce faithful local abstractions rather than scalar placeholder aliases.

\begin{theorem}[Physics formalization target]
\label{thm:physics:phyx_mini_0291:target}
\lean{PhyXMiniProblems.ProblemPhyXMini0291.problem_phyx_mini_0291}
\uses{def:physics:phyx-mini-0291:phyxminiproblems-problemphyxmini0291-timeinseconds, def:physics:phyx-mini-0291:phyxminiproblems-problemphyxmini0291-travelingstringwavesetup, def:physics:phyx-mini-0291:phyxminiproblems-problemphyxmini0291-matchestravelingwavescenario, def:physics:phyx-mini-0291:phyxminiproblems-problemphyxmini0291-matchesproblemreadouts, def:physics:phyx-mini-0291:phyxminiproblems-problemphyxmini0291-matchessuppliedfigure, def:physics:phyx-mini-0291:phyxminiproblems-problemphyxmini0291-hasphysicalstringwaveparameters, def:physics:phyx-mini-0291:phyxminiproblems-problemphyxmini0291-satisfiesuniformstringwavelaws, def:physics:phyx-mini-0291:phyxminiproblems-problemphyxmini0291-answerchoice, def:physics:phyx-mini-0291:phyxminiproblems-problemphyxmini0291-isnearestdisplayedperiod}
The assigned autoformalize agent should translate this physics problem into Lean declarations in the covered file, with theorem and lemma proof bodies written as `by sorry`.
\end{theorem}
\begin{proof}
This is an autoformalization task, not a proof task. Produce statements that can later be proved by the physics prover without weakening the physical model.
\end{proof}

% --- Archon declaration topology begin ---
\section{Declaration topology}
\begin{definition}[Length Quantity]
  \label{def:physics:phyx-mini-0291:phyxminiproblems-problemphyxmini0291-lengthquantity}
  \lean{PhyXMiniProblems.ProblemPhyXMini0291.LengthQuantity}
  A nonnegative physical length, independent of the unit used to read it
\end{definition}
\begin{proof}
  This is the declared model definition of Length Quantity.
\end{proof}
\begin{definition}[Signed Length Quantity]
  \label{def:physics:phyx-mini-0291:phyxminiproblems-problemphyxmini0291-signedlengthquantity}
  \lean{PhyXMiniProblems.ProblemPhyXMini0291.SignedLengthQuantity}
  A signed one-dimensional position or transverse displacement
\end{definition}
\begin{proof}
  This is the declared model definition of Signed Length Quantity.
\end{proof}
\begin{definition}[Time Quantity]
  \label{def:physics:phyx-mini-0291:phyxminiproblems-problemphyxmini0291-timequantity}
  \lean{PhyXMiniProblems.ProblemPhyXMini0291.TimeQuantity}
  A nonnegative physical duration
\end{definition}
\begin{proof}
  This is the declared model definition of Time Quantity.
\end{proof}
\begin{definition}[Linear Mass Density Quantity]
  \label{def:physics:phyx-mini-0291:phyxminiproblems-problemphyxmini0291-linearmassdensityquantity}
  \lean{PhyXMiniProblems.ProblemPhyXMini0291.LinearMassDensityQuantity}
  Mass per unit length of the uniform string
\end{definition}
\begin{proof}
  This is the declared model definition of Linear Mass Density Quantity.
\end{proof}
\begin{definition}[Tension Quantity]
  \label{def:physics:phyx-mini-0291:phyxminiproblems-problemphyxmini0291-tensionquantity}
  \lean{PhyXMiniProblems.ProblemPhyXMini0291.TensionQuantity}
  A tensile force, with dimension mass * length / time\textasciicircum{}2
\end{definition}
\begin{proof}
  This is the declared model definition of Tension Quantity.
\end{proof}
\begin{definition}[Speed Quantity]
  \label{def:physics:phyx-mini-0291:phyxminiproblems-problemphyxmini0291-speedquantity}
  \lean{PhyXMiniProblems.ProblemPhyXMini0291.SpeedQuantity}
  The nonnegative magnitude of the transverse-wave propagation velocity
\end{definition}
\begin{proof}
  This is the declared model definition of Speed Quantity.
\end{proof}
\begin{definition}[length Readout]
  \label{def:physics:phyx-mini-0291:phyxminiproblems-problemphyxmini0291-lengthreadout}
  \lean{PhyXMiniProblems.ProblemPhyXMini0291.lengthReadout}
  \uses{def:physics:phyx-mini-0291:phyxminiproblems-problemphyxmini0291-lengthquantity}
  Read a physical length in a selected length unit
\end{definition}
\begin{proof}
  This is the declared model definition of length Readout.
\end{proof}
\begin{definition}[signed Length Readout]
  \label{def:physics:phyx-mini-0291:phyxminiproblems-problemphyxmini0291-signedlengthreadout}
  \lean{PhyXMiniProblems.ProblemPhyXMini0291.signedLengthReadout}
  \uses{def:physics:phyx-mini-0291:phyxminiproblems-problemphyxmini0291-signedlengthquantity}
  Read a signed position or displacement in a selected length unit
\end{definition}
\begin{proof}
  This is the declared model definition of signed Length Readout.
\end{proof}
\begin{definition}[time Readout]
  \label{def:physics:phyx-mini-0291:phyxminiproblems-problemphyxmini0291-timereadout}
  \lean{PhyXMiniProblems.ProblemPhyXMini0291.timeReadout}
  \uses{def:physics:phyx-mini-0291:phyxminiproblems-problemphyxmini0291-timequantity}
  Read a physical duration in a selected time unit
\end{definition}
\begin{proof}
  This is the declared model definition of time Readout.
\end{proof}
\begin{definition}[linear Mass Density Readout]
  \label{def:physics:phyx-mini-0291:phyxminiproblems-problemphyxmini0291-linearmassdensityreadout}
  \lean{PhyXMiniProblems.ProblemPhyXMini0291.linearMassDensityReadout}
  \uses{def:physics:phyx-mini-0291:phyxminiproblems-problemphyxmini0291-linearmassdensityquantity}
  Read linear mass density in selected mass-per-length units
\end{definition}
\begin{proof}
  This is the declared model definition of linear Mass Density Readout.
\end{proof}
\begin{definition}[tension Readout]
  \label{def:physics:phyx-mini-0291:phyxminiproblems-problemphyxmini0291-tensionreadout}
  \lean{PhyXMiniProblems.ProblemPhyXMini0291.tensionReadout}
  \uses{def:physics:phyx-mini-0291:phyxminiproblems-problemphyxmini0291-tensionquantity}
  Read tension in the coherent force unit induced by the base units
\end{definition}
\begin{proof}
  This is the declared model definition of tension Readout.
\end{proof}
\begin{definition}[speed Readout]
  \label{def:physics:phyx-mini-0291:phyxminiproblems-problemphyxmini0291-speedreadout}
  \lean{PhyXMiniProblems.ProblemPhyXMini0291.speedReadout}
  \uses{def:physics:phyx-mini-0291:phyxminiproblems-problemphyxmini0291-speedquantity}
  Read wave speed in a selected length unit per selected time unit
\end{definition}
\begin{proof}
  This is the declared model definition of speed Readout.
\end{proof}
\begin{definition}[length In Meters]
  \label{def:physics:phyx-mini-0291:phyxminiproblems-problemphyxmini0291-lengthinmeters}
  \lean{PhyXMiniProblems.ProblemPhyXMini0291.lengthInMeters}
  \uses{def:physics:phyx-mini-0291:phyxminiproblems-problemphyxmini0291-lengthquantity, def:physics:phyx-mini-0291:phyxminiproblems-problemphyxmini0291-lengthreadout}
  Meter readout of a physical length
\end{definition}
\begin{proof}
  This is the declared model definition of length In Meters.
\end{proof}
\begin{definition}[length In Centimeters]
  \label{def:physics:phyx-mini-0291:phyxminiproblems-problemphyxmini0291-lengthincentimeters}
  \lean{PhyXMiniProblems.ProblemPhyXMini0291.lengthInCentimeters}
  \uses{def:physics:phyx-mini-0291:phyxminiproblems-problemphyxmini0291-lengthquantity, def:physics:phyx-mini-0291:phyxminiproblems-problemphyxmini0291-lengthreadout}
  Centimeter readout used by both axes of the supplied graph
\end{definition}
\begin{proof}
  This is the declared model definition of length In Centimeters.
\end{proof}
\begin{definition}[time In Seconds]
  \label{def:physics:phyx-mini-0291:phyxminiproblems-problemphyxmini0291-timeinseconds}
  \lean{PhyXMiniProblems.ProblemPhyXMini0291.timeInSeconds}
  \uses{def:physics:phyx-mini-0291:phyxminiproblems-problemphyxmini0291-timequantity, def:physics:phyx-mini-0291:phyxminiproblems-problemphyxmini0291-timereadout}
  Second readout of a physical duration
\end{definition}
\begin{proof}
  This is the declared model definition of time In Seconds.
\end{proof}
\begin{definition}[linear Mass Density In Grams Per Meter]
  \label{def:physics:phyx-mini-0291:phyxminiproblems-problemphyxmini0291-linearmassdensityingramspermeter}
  \lean{PhyXMiniProblems.ProblemPhyXMini0291.linearMassDensityInGramsPerMeter}
  \uses{def:physics:phyx-mini-0291:phyxminiproblems-problemphyxmini0291-linearmassdensityquantity, def:physics:phyx-mini-0291:phyxminiproblems-problemphyxmini0291-linearmassdensityreadout}
  Linear-density readout in grams per meter, as printed in the problem
\end{definition}
\begin{proof}
  This is the declared model definition of linear Mass Density In Grams Per Meter.
\end{proof}
\begin{definition}[linear Mass Density In Kilograms Per Meter]
  \label{def:physics:phyx-mini-0291:phyxminiproblems-problemphyxmini0291-linearmassdensityinkilogramspermeter}
  \lean{PhyXMiniProblems.ProblemPhyXMini0291.linearMassDensityInKilogramsPerMeter}
  \uses{def:physics:phyx-mini-0291:phyxminiproblems-problemphyxmini0291-linearmassdensityquantity, def:physics:phyx-mini-0291:phyxminiproblems-problemphyxmini0291-linearmassdensityreadout}
  SI linear-density readout in kilograms per meter
\end{definition}
\begin{proof}
  This is the declared model definition of linear Mass Density In Kilograms Per Meter.
\end{proof}
\begin{definition}[tension In Newtons]
  \label{def:physics:phyx-mini-0291:phyxminiproblems-problemphyxmini0291-tensioninnewtons}
  \lean{PhyXMiniProblems.ProblemPhyXMini0291.tensionInNewtons}
  \uses{def:physics:phyx-mini-0291:phyxminiproblems-problemphyxmini0291-tensionquantity, def:physics:phyx-mini-0291:phyxminiproblems-problemphyxmini0291-tensionreadout}
  SI tension readout in newtons
\end{definition}
\begin{proof}
  This is the declared model definition of tension In Newtons.
\end{proof}
\begin{definition}[speed In Meters Per Second]
  \label{def:physics:phyx-mini-0291:phyxminiproblems-problemphyxmini0291-speedinmeterspersecond}
  \lean{PhyXMiniProblems.ProblemPhyXMini0291.speedInMetersPerSecond}
  \uses{def:physics:phyx-mini-0291:phyxminiproblems-problemphyxmini0291-speedquantity, def:physics:phyx-mini-0291:phyxminiproblems-problemphyxmini0291-speedreadout}
  SI speed readout in meters per second
\end{definition}
\begin{proof}
  This is the declared model definition of speed In Meters Per Second.
\end{proof}
\begin{definition}[Wave Medium Kind]
  \label{def:physics:phyx-mini-0291:phyxminiproblems-problemphyxmini0291-wavemediumkind}
  \lean{PhyXMiniProblems.ProblemPhyXMini0291.WaveMediumKind}
  The material object supporting the wave
\end{definition}
\begin{proof}
  This is the declared model definition of Wave Medium Kind.
\end{proof}
\begin{definition}[Wave Motion Kind]
  \label{def:physics:phyx-mini-0291:phyxminiproblems-problemphyxmini0291-wavemotionkind}
  \lean{PhyXMiniProblems.ProblemPhyXMini0291.WaveMotionKind}
  The stated waveform and polarization
\end{definition}
\begin{proof}
  This is the declared model definition of Wave Motion Kind.
\end{proof}
\begin{definition}[Propagation Direction]
  \label{def:physics:phyx-mini-0291:phyxminiproblems-problemphyxmini0291-propagationdirection}
  \lean{PhyXMiniProblems.ProblemPhyXMini0291.PropagationDirection}
  Direction of travel along the horizontal coordinate axis
\end{definition}
\begin{proof}
  This is the declared model definition of Propagation Direction.
\end{proof}
\begin{definition}[Graph Axis Symbol]
  \label{def:physics:phyx-mini-0291:phyxminiproblems-problemphyxmini0291-graphaxissymbol}
  \lean{PhyXMiniProblems.ProblemPhyXMini0291.GraphAxisSymbol}
  Symbols printed beside the two axes in the primary raster
\end{definition}
\begin{proof}
  This is the declared model definition of Graph Axis Symbol.
\end{proof}
\begin{definition}[Figure Label]
  \label{def:physics:phyx-mini-0291:phyxminiproblems-problemphyxmini0291-figurelabel}
  \lean{PhyXMiniProblems.ProblemPhyXMini0291.FigureLabel}
  Text labels visible in the supplied graph
\end{definition}
\begin{proof}
  This is the declared model definition of Figure Label.
\end{proof}
\begin{definition}[Graph Point Label]
  \label{def:physics:phyx-mini-0291:phyxminiproblems-problemphyxmini0291-graphpointlabel}
  \lean{PhyXMiniProblems.ProblemPhyXMini0291.GraphPointLabel}
  Grid-readable points spanning one complete spatial cycle at t = 0
\end{definition}
\begin{proof}
  This is the declared model definition of Graph Point Label.
\end{proof}
\begin{definition}[Displacement Curve Shape]
  \label{def:physics:phyx-mini-0291:phyxminiproblems-problemphyxmini0291-displacementcurveshape}
  \lean{PhyXMiniProblems.ProblemPhyXMini0291.DisplacementCurveShape}
  Qualitative shape of the displacement-versus-position trace
\end{definition}
\begin{proof}
  This is the declared model definition of Displacement Curve Shape.
\end{proof}
\begin{definition}[displayed Graph Coordinate]
  \label{def:physics:phyx-mini-0291:phyxminiproblems-problemphyxmini0291-displayedgraphcoordinate}
  \lean{PhyXMiniProblems.ProblemPhyXMini0291.displayedGraphCoordinate}
  \uses{def:physics:phyx-mini-0291:phyxminiproblems-problemphyxmini0291-graphpointlabel}
  Idealized (x,y) coordinates read from the primary grid, in centimeters. The two crests are 40 cm apart. Since y\_s = 4 cm marks the horizontal grid line one division below the crest, the plotted amplitude is 5 cm
\end{definition}
\begin{proof}
  This is the declared model definition of displayed Graph Coordinate.
\end{proof}
\begin{definition}[String Wave Figure]
  \label{def:physics:phyx-mini-0291:phyxminiproblems-problemphyxmini0291-stringwavefigure}
  \lean{PhyXMiniProblems.ProblemPhyXMini0291.StringWaveFigure}
  \uses{def:physics:phyx-mini-0291:phyxminiproblems-problemphyxmini0291-lengthquantity, def:physics:phyx-mini-0291:phyxminiproblems-problemphyxmini0291-signedlengthquantity, def:physics:phyx-mini-0291:phyxminiproblems-problemphyxmini0291-timequantity, def:physics:phyx-mini-0291:phyxminiproblems-problemphyxmini0291-graphaxissymbol, def:physics:phyx-mini-0291:phyxminiproblems-problemphyxmini0291-figurelabel, def:physics:phyx-mini-0291:phyxminiproblems-problemphyxmini0291-graphpointlabel, def:physics:phyx-mini-0291:phyxminiproblems-problemphyxmini0291-displacementcurveshape}
  Dimensionful axes, marked grid points, and metadata of the primary graph
\end{definition}
\begin{proof}
  This is the declared model definition of String Wave Figure.
\end{proof}
\begin{definition}[Traveling String Wave Setup]
  \label{def:physics:phyx-mini-0291:phyxminiproblems-problemphyxmini0291-travelingstringwavesetup}
  \lean{PhyXMiniProblems.ProblemPhyXMini0291.TravelingStringWaveSetup}
  \uses{def:physics:phyx-mini-0291:phyxminiproblems-problemphyxmini0291-lengthquantity, def:physics:phyx-mini-0291:phyxminiproblems-problemphyxmini0291-signedlengthquantity, def:physics:phyx-mini-0291:phyxminiproblems-problemphyxmini0291-timequantity, def:physics:phyx-mini-0291:phyxminiproblems-problemphyxmini0291-linearmassdensityquantity, def:physics:phyx-mini-0291:phyxminiproblems-problemphyxmini0291-tensionquantity, def:physics:phyx-mini-0291:phyxminiproblems-problemphyxmini0291-speedquantity, def:physics:phyx-mini-0291:phyxminiproblems-problemphyxmini0291-wavemediumkind, def:physics:phyx-mini-0291:phyxminiproblems-problemphyxmini0291-wavemotionkind, def:physics:phyx-mini-0291:phyxminiproblems-problemphyxmini0291-propagationdirection, def:physics:phyx-mini-0291:phyxminiproblems-problemphyxmini0291-stringwavefigure}
  The independent physical observables of the traveling wave. In particular, the requested period is not defined from an answer choice or target value
\end{definition}
\begin{proof}
  This is the declared model definition of Traveling String Wave Setup.
\end{proof}
\begin{definition}[Matches Traveling Wave Scenario]
  \label{def:physics:phyx-mini-0291:phyxminiproblems-problemphyxmini0291-matchestravelingwavescenario}
  \lean{PhyXMiniProblems.ProblemPhyXMini0291.MatchesTravelingWaveScenario}
  \uses{def:physics:phyx-mini-0291:phyxminiproblems-problemphyxmini0291-travelingstringwavesetup}
  Categorical facts stated in the physical scenario
\end{definition}
\begin{proof}
  This is the declared model definition of Matches Traveling Wave Scenario.
\end{proof}
\begin{definition}[Matches Problem Readouts]
  \label{def:physics:phyx-mini-0291:phyxminiproblems-problemphyxmini0291-matchesproblemreadouts}
  \lean{PhyXMiniProblems.ProblemPhyXMini0291.MatchesProblemReadouts}
  \uses{def:physics:phyx-mini-0291:phyxminiproblems-problemphyxmini0291-lengthincentimeters, def:physics:phyx-mini-0291:phyxminiproblems-problemphyxmini0291-linearmassdensityingramspermeter, def:physics:phyx-mini-0291:phyxminiproblems-problemphyxmini0291-tensioninnewtons, def:physics:phyx-mini-0291:phyxminiproblems-problemphyxmini0291-travelingstringwavesetup}
  Numerical readouts supplied by the prose. No period, speed, or wavelength value is present here
\end{definition}
\begin{proof}
  This is the declared model definition of Matches Problem Readouts.
\end{proof}
\begin{definition}[Matches Supplied Figure]
  \label{def:physics:phyx-mini-0291:phyxminiproblems-problemphyxmini0291-matchessuppliedfigure}
  \lean{PhyXMiniProblems.ProblemPhyXMini0291.MatchesSuppliedFigure}
  \uses{def:physics:phyx-mini-0291:phyxminiproblems-problemphyxmini0291-lengthreadout, def:physics:phyx-mini-0291:phyxminiproblems-problemphyxmini0291-signedlengthreadout, def:physics:phyx-mini-0291:phyxminiproblems-problemphyxmini0291-lengthincentimeters, def:physics:phyx-mini-0291:phyxminiproblems-problemphyxmini0291-timeinseconds, def:physics:phyx-mini-0291:phyxminiproblems-problemphyxmini0291-figurelabel, def:physics:phyx-mini-0291:phyxminiproblems-problemphyxmini0291-graphpointlabel, def:physics:phyx-mini-0291:phyxminiproblems-problemphyxmini0291-displayedgraphcoordinate, def:physics:phyx-mini-0291:phyxminiproblems-problemphyxmini0291-travelingstringwavesetup}
  Primary-image evidence. The wavelength relation is a geometric readout from the distinguished pair of successive crests, not a premise about the period
\end{definition}
\begin{proof}
  This is the declared model definition of Matches Supplied Figure.
\end{proof}
\begin{definition}[Has Physical String Wave Parameters]
  \label{def:physics:phyx-mini-0291:phyxminiproblems-problemphyxmini0291-hasphysicalstringwaveparameters}
  \lean{PhyXMiniProblems.ProblemPhyXMini0291.HasPhysicalStringWaveParameters}
  \uses{def:physics:phyx-mini-0291:phyxminiproblems-problemphyxmini0291-lengthinmeters, def:physics:phyx-mini-0291:phyxminiproblems-problemphyxmini0291-timeinseconds, def:physics:phyx-mini-0291:phyxminiproblems-problemphyxmini0291-linearmassdensityinkilogramspermeter, def:physics:phyx-mini-0291:phyxminiproblems-problemphyxmini0291-tensioninnewtons, def:physics:phyx-mini-0291:phyxminiproblems-problemphyxmini0291-speedinmeterspersecond, def:physics:phyx-mini-0291:phyxminiproblems-problemphyxmini0291-travelingstringwavesetup}
  Positivity and nondegeneracy of the physical wave parameters
\end{definition}
\begin{proof}
  This is the declared model definition of Has Physical String Wave Parameters.
\end{proof}
\begin{definition}[Satisfies Uniform String Wave Laws]
  \label{def:physics:phyx-mini-0291:phyxminiproblems-problemphyxmini0291-satisfiesuniformstringwavelaws}
  \lean{PhyXMiniProblems.ProblemPhyXMini0291.SatisfiesUniformStringWaveLaws}
  \uses{def:physics:phyx-mini-0291:phyxminiproblems-problemphyxmini0291-lengthreadout, def:physics:phyx-mini-0291:phyxminiproblems-problemphyxmini0291-timereadout, def:physics:phyx-mini-0291:phyxminiproblems-problemphyxmini0291-linearmassdensityreadout, def:physics:phyx-mini-0291:phyxminiproblems-problemphyxmini0291-tensionreadout, def:physics:phyx-mini-0291:phyxminiproblems-problemphyxmini0291-speedreadout, def:physics:phyx-mini-0291:phyxminiproblems-problemphyxmini0291-travelingstringwavesetup}
  The two governing laws used by the textbook solution, in arbitrary coherent units: T = μ v² and λ = v P. These general relations do not assign any numerical period or answer label
\end{definition}
\begin{proof}
  This is the declared model definition of Satisfies Uniform String Wave Laws.
\end{proof}
\begin{lemma}[wavelength In Meters eq two fifths]
  \label{lem:physics:phyx-mini-0291:phyxminiproblems-problemphyxmini0291-wavelengthinmeters-eq-two-fifths}
  \lean{PhyXMiniProblems.ProblemPhyXMini0291.wavelengthInMeters_eq_two_fifths}
  \uses{def:physics:phyx-mini-0291:phyxminiproblems-problemphyxmini0291-lengthinmeters, def:physics:phyx-mini-0291:phyxminiproblems-problemphyxmini0291-travelingstringwavesetup, def:physics:phyx-mini-0291:phyxminiproblems-problemphyxmini0291-matchessuppliedfigure}
  Successive crests at 5 cm and 45 cm give wavelength 2/5 m
\end{lemma}
\begin{proof}
  The conclusion follows from the governing relations and figure readouts named in the statement of wavelength In Meters eq two fifths.
\end{proof}
\begin{lemma}[propagation Speed In Meters Per Second eq twelve]
  \label{lem:physics:phyx-mini-0291:phyxminiproblems-problemphyxmini0291-propagationspeedinmeterspersecond-eq-twelve}
  \lean{PhyXMiniProblems.ProblemPhyXMini0291.propagationSpeedInMetersPerSecond_eq_twelve}
  \uses{def:physics:phyx-mini-0291:phyxminiproblems-problemphyxmini0291-speedinmeterspersecond, def:physics:phyx-mini-0291:phyxminiproblems-problemphyxmini0291-travelingstringwavesetup, def:physics:phyx-mini-0291:phyxminiproblems-problemphyxmini0291-matchesproblemreadouts, def:physics:phyx-mini-0291:phyxminiproblems-problemphyxmini0291-hasphysicalstringwaveparameters, def:physics:phyx-mini-0291:phyxminiproblems-problemphyxmini0291-satisfiesuniformstringwavelaws}
  The string-wave law and prose readouts give propagation speed 12 m/s
\end{lemma}
\begin{proof}
  The conclusion follows from the governing relations and figure readouts named in the statement of propagation Speed In Meters Per Second eq twelve.
\end{proof}
\begin{lemma}[wave Period In Seconds eq one thirtieth]
  \label{lem:physics:phyx-mini-0291:phyxminiproblems-problemphyxmini0291-waveperiodinseconds-eq-one-thirtieth}
  \lean{PhyXMiniProblems.ProblemPhyXMini0291.wavePeriodInSeconds_eq_one_thirtieth}
  \uses{def:physics:phyx-mini-0291:phyxminiproblems-problemphyxmini0291-timeinseconds, def:physics:phyx-mini-0291:phyxminiproblems-problemphyxmini0291-travelingstringwavesetup, def:physics:phyx-mini-0291:phyxminiproblems-problemphyxmini0291-matchesproblemreadouts, def:physics:phyx-mini-0291:phyxminiproblems-problemphyxmini0291-matchessuppliedfigure, def:physics:phyx-mini-0291:phyxminiproblems-problemphyxmini0291-hasphysicalstringwaveparameters, def:physics:phyx-mini-0291:phyxminiproblems-problemphyxmini0291-satisfiesuniformstringwavelaws}
  Combining λ = 2/5 m, v = 12 m/s, and λ = v P gives P = 1/30 s
\end{lemma}
\begin{proof}
  The conclusion follows from the governing relations and figure readouts named in the statement of wave Period In Seconds eq one thirtieth.
\end{proof}
\begin{definition}[Answer Choice]
  \label{def:physics:phyx-mini-0291:phyxminiproblems-problemphyxmini0291-answerchoice}
  \lean{PhyXMiniProblems.ProblemPhyXMini0291.AnswerChoice}
  The four answer labels printed with the problem
\end{definition}
\begin{proof}
  This is the declared model definition of Answer Choice.
\end{proof}
\begin{definition}[displayed Period In Seconds]
  \label{def:physics:phyx-mini-0291:phyxminiproblems-problemphyxmini0291-displayedperiodinseconds}
  \lean{PhyXMiniProblems.ProblemPhyXMini0291.displayedPeriodInSeconds}
  \uses{def:physics:phyx-mini-0291:phyxminiproblems-problemphyxmini0291-answerchoice}
  Displayed period values, in seconds
\end{definition}
\begin{proof}
  This is the declared model definition of displayed Period In Seconds.
\end{proof}
\begin{definition}[recorded Dataset Answer]
  \label{def:physics:phyx-mini-0291:phyxminiproblems-problemphyxmini0291-recordeddatasetanswer}
  \lean{PhyXMiniProblems.ProblemPhyXMini0291.recordedDatasetAnswer}
  \uses{def:physics:phyx-mini-0291:phyxminiproblems-problemphyxmini0291-answerchoice}
  The answer label recorded by the source dataset; this is not a premise
\end{definition}
\begin{proof}
  This is the declared model definition of recorded Dataset Answer.
\end{proof}
\begin{definition}[Is Nearest Displayed Period]
  \label{def:physics:phyx-mini-0291:phyxminiproblems-problemphyxmini0291-isnearestdisplayedperiod}
  \lean{PhyXMiniProblems.ProblemPhyXMini0291.IsNearestDisplayedPeriod}
  \uses{def:physics:phyx-mini-0291:phyxminiproblems-problemphyxmini0291-timequantity, def:physics:phyx-mini-0291:phyxminiproblems-problemphyxmini0291-timeinseconds, def:physics:phyx-mini-0291:phyxminiproblems-problemphyxmini0291-answerchoice, def:physics:phyx-mini-0291:phyxminiproblems-problemphyxmini0291-displayedperiodinseconds}
  A displayed answer is at least as close as every other displayed value
\end{definition}
\begin{proof}
  This is the declared model definition of Is Nearest Displayed Period.
\end{proof}
% --- Archon declaration topology end ---
% --- Archon physics formalization source end ---
